\subsection*{Notes for IX. EXPERIMENTAL PROTOCOL}
\paragraph{Objective}
Specify the full experimental design and auditing details: benchmarks, baselines, budgets, seeds, metrics, statistical plan, diagnostics, and reproducibility artifacts.

\paragraph{Critical clarity to add (highest priority)}
The protocol must not allow the method to be interpreted as “P3 + Pareto filtering”. It must make explicit:
\begin{itemize}
  \item which NSGA-II mechanisms you retain (population/archive semantics, nondominated sorting, crowding, replacement),
  \item what P3 replaces/augments (variation operator and proposal distribution),
  \item how multi-fidelity affects elite definition and model updates.
\end{itemize}

\subsection*{Notes for IX: Protocol at a glance}
\paragraph{Write next}
\begin{itemize}
  \item One compact bullet list: benchmarks, budgets, \#runs, proposed method, baselines, primary metric(s), stats tests, artifact.
  \item One fairness bullet list: identical evaluator, encoding, ASHA schedule, budget accounting, parallelism, initialization.
\end{itemize}

\subsection*{Notes for IX: Notation}
\paragraph{Write next}
\begin{itemize}
  \item Define symbols used throughout: $\Lambda$, $x$, $D(x)$, $\mathbf{f}(x)$, $B$, $R$, $W$, $A$, $\rho$, $U$, rungs, archive $\mathcal{A}$, elite set $\mathcal{E}$, distribution $q_\theta$.
  \item Include a short “note on "units”": what counts as one evaluation/cost unit.
\end{itemize}

\subsection*{Notes for IX.A Empirical Methodology and Study Type}
\paragraph{Write next}
\begin{itemize}
  \item Controlled benchmarking of stochastic optimizers with repeated independent runs.
  \item Pre-specified hypotheses and falsification criteria (short form).
  \item Seed policy: paired runs/common random numbers, evaluator determinism controls.
\end{itemize}

\subsection*{Notes for IX.C Algorithms Compared (Proposed + Baselines)}
\paragraph{Write next}
\begin{itemize}
  \item Proposed: P3Net definition in 5--8 lines, referencing the full Method section.
  \item Baseline list with feasibility notes:
  \begin{itemize}
    \item Random Search (RS), RS+ASHA,
    \item NSGA-II/NSGA-Net-style baseline on the same encoding/evaluator,
    \item P3 alone (without NSGA-II multiobjective selection mechanics),
    \item BOHB/TPE/SMAC only if you can provide concrete encoding mapping and conditional support.
  \end{itemize}
  \item For each baseline: code source/version, key settings, tuning policy (none vs bounded and equal across methods).
\end{itemize}

\subsection*{Notes for IX.B Benchmarks / Problem Instances}
\paragraph{Write next}
\begin{itemize}
  \item Enumerate benchmark tasks (e.g., CIFAR-10/100) and evaluation pipeline.
  \item Define ASHA schedule: rung epochs, promotion rules, max epochs.
  \item Define the search space precisely and reference the formal definition in the formulation section.
\end{itemize}


\paragraph{Fairness plan (must be explicit)}
\begin{itemize}
  \item Same encoding and constraints handling (or explicitly different with justification).
  \item Same budget accounting, including failures and resampling overhead.
  \item The same ASHA schedule and parallelism policy where applicable.
\end{itemize}

\subsection*{Notes for IX.D Experimental Design (Budgets, Seeds, Repeats, Stopping)}
\paragraph{Write next}
\begin{itemize}
  \item Budgets $B$ (list values) and what they count.
  \item Runs $R$ (e.g., 10) and seed list; paired design where possible.
  \item Initialization $n_0$ and how it is generated (valid random architectures).
  \item Stopping rules, failure handling, and how incomplete trials are counted.
\end{itemize}

\subsection*{Notes for IX.E Evaluation Metrics and Comparison Targets}
\paragraph{Write next}
\begin{itemize}
  \item Primary metric(s):
  \begin{itemize}
    \item Single-objective: best-sofar regret vs evaluations; fixed-budget regret.
    \item Multiobjective: HV or HV gap vs evaluations; fixed-budget HV gap.
  \end{itemize}
  \item Secondary metrics: $\epsilon$-indicator, coverage, time-to-target (with censoring rules), failure rate.
  \item Define the HV reference point and how it is chosen.
\end{itemize}

\subsection*{Notes for IX.F Statistical Analysis Plan (Pre-specified)}
\paragraph{Write next}
\begin{itemize}
  \item Paired Wilcoxon signed-rank tests for summary metrics per-run.
  \item Holm--Bonferroni correction across multiple comparisons.
  \item Effect sizes (e.g., Cliff’s delta) and uncertainty reporting (IQR/CI).
  \item Time-to-target: censoring/survival-style handling if used.
  \item Failure handling: define how OOM/divergence is included in metrics/tests.
\end{itemize}

\subsection*{Notes for IX.G Error Analysis / Diagnostic Plots}
\paragraph{Write next}
\begin{itemize}
  \item Diversity/entropy of sampled tokens over time; duplicate rate; archive turnover.
  \item Invalid sample rate (pre/post repair) and resampling attempts distribution.
  \item ASHA rung distribution per method: how many reach the top rung (supports the elite learning rule).
  \item Failure taxonomy: OOM, divergence, invalid genotype, training instability.
\end{itemize}

\subsection*{Notes for IX.H Reproducibility Statement + Artifact Plan}
\paragraph{Write next}
\begin{itemize}
  \item Exact reproducibility: pinned environment + seeds $\rightarrow$ regenerate all figures/tables.
  \item Broad reproducibility: qualitative trends under different hardware (state expected variance).
  \item Artifact Contents
  \begin{itemize}
    \item code for encoding/decoding/repair + P3Net,
    \item baseline wrappers and configurations,
    \item raw per-evaluation logs (architecture hash, rung, objectives, seed, failure codes),
    \item scripts to regenerate each figure/table,
    \item statistics scripts (tests + corrections + effect sizes),
    \item container/lockfile with pinned versions.
  \end{itemize}
\end{itemize}
